\documentclass[10pt,a4paper]{article}

% --- Кодировка, шрифты и русский язык ---
\usepackage{cmap}
\usepackage[T2A]{fontenc}
\usepackage[utf8]{inputenc}
\usepackage[english,russian]{babel}
\usepackage{microtype}

% --- Математика ---
\usepackage{amsmath,amssymb,amsfonts,mathtools}

% --- Геометрия и оформление ---
\usepackage[a4paper,left=15mm,right=15mm,top=18mm,bottom=18mm,headheight=14pt]{geometry}
\usepackage[dvipsnames,svgnames]{xcolor}
\usepackage{enumitem}
\usepackage{booktabs}
\usepackage{array}
\usepackage{tabularx}
\usepackage{fancyhdr}
\usepackage{titlesec}
\usepackage[most]{tcolorbox}
\usepackage[colorlinks=true,linkcolor=NavyBlue,urlcolor=NavyBlue,citecolor=NavyBlue]{hyperref}

% --- Цвета и блоки ---
\definecolor{AccentBlue}{HTML}{1B365D}
\definecolor{SubtleBg}{HTML}{F4F6F9}
\definecolor{BorderBlue}{HTML}{B8C7DC}

\newtcolorbox{summarybox}[1][]{
  enhanced,
  breakable,
  colback=SubtleBg,
  colframe=BorderBlue,
  boxrule=0.6pt,
  arc=2pt,
  left=6pt,right=6pt,top=5pt,bottom=5pt,
  fonttitle=\bfseries\color{AccentBlue},
  coltitle=AccentBlue,
  colbacktitle=SubtleBg,
  attach title to upper={.\ },
  #1
}

% --- Настройка заголовков и списков ---
\titleformat{\section}
  {\Large\bfseries\color{AccentBlue}}
  {\thesection.}{0.5em}{}[\vspace{2pt}\titlerule]
\titlespacing*{\section}{0pt}{12pt}{6pt}

\titleformat{\subsection}
  {\large\bfseries\color{AccentBlue!90!black}}
  {\thesubsection.}{0.5em}{}
\titlespacing*{\subsection}{0pt}{9pt}{4pt}

\titleformat{\subsubsection}
  {\normalsize\bfseries}
  {\thesubsubsection.}{0.5em}{}
\titlespacing*{\subsubsection}{0pt}{6pt}{2pt}

\setlist[itemize]{leftmargin=1.5em,itemsep=2pt,topsep=3pt,parsep=0pt,beginpenalty=10000}
\setlist[enumerate]{leftmargin=1.7em,itemsep=2pt,topsep=3pt,parsep=0pt,beginpenalty=10000}
\setlength{\parindent}{0pt}
\setlength{\parskip}{3pt}

% --- Колонтитулы ---
\pagestyle{fancy}
\fancyhf{}
\fancyhead[L]{\footnotesize\textcolor{gray}{TeamBook: Тервер и комбинаторика}}
\fancyhead[R]{\footnotesize\textcolor{gray}{\nouppercase{\leftmark}}}
\fancyfoot[C]{\small\thepage}
\renewcommand{\headrulewidth}{0.4pt}

% --- Математические макросы и русская типографика ---
\renewcommand{\le}{\leqslant}
\renewcommand{\ge}{\geqslant}
\renewcommand{\emptyset}{\varnothing}

\newcommand{\N}{\mathbb{N}}
\newcommand{\Z}{\mathbb{Z}}
\newcommand{\R}{\mathbb{R}}
\newcommand{\F}{\mathbb{F}}
\newcommand{\K}{\mathbb{K}}
\newcommand{\E}{\mathbb{E}}
\newcommand{\Var}{\operatorname{Var}}
\newcommand{\Ind}{\mathbb{I}}
\newcommand{\iverson}[1]{[\,#1\,]}
\newcommand{\fall}[2]{#1^{\underline{#2}}}
\newcommand{\rise}[2]{#1^{\overline{#2}}}
\newcommand{\stirlingfirst}[2]{\genfrac{[}{]}{0pt}{}{#1}{#2}}
\newcommand{\stirlingsecond}[2]{\genfrac{\{}{\}}{0pt}{}{#1}{#2}}
\DeclareMathOperator{\lcm}{lcm}
\DeclareMathOperator{\rk}{rank}
\DeclareMathOperator{\unrk}{unrank}
\DeclareMathOperator{\inv}{inv}
\DeclareMathOperator{\Fix}{Fix}
\DeclareMathOperator{\Stab}{Stab}
\DeclareMathOperator{\sgn}{sgn}
\DeclareMathOperator{\diag}{diag}
\DeclareMathOperator{\popcnt}{popcount}
\DeclareMathOperator{\Geom}{Geom}
\DeclareMathOperator{\Unif}{Unif}

\begin{document}
\thispagestyle{plain}

% --- Заголовок ---
\begin{center}
  {\LARGE\bfseries\color{AccentBlue} Теория вероятностей и комбинаторика}\\[4pt]
  {\normalsize Краткий справочник по теории для олимпиадной информатики (TeamBook, CF 1300--2300+)}
\end{center}
\vspace{2pt}

\begin{summarybox}[title={Базовые обозначения и соглашения}]
\begin{itemize}
  \item $\N = \{1, 2, \dots\}$, $\N_0 = \{0, 1, 2, \dots\}$, $\Z$ — целые числа, $\R$ — вещественные числа; $[n] = \{1, 2, \dots, n\}$ при $n \ge 1$ и $[0] = \emptyset$.
  \item $|A|$ — мощность конечного множества $A$; $2^A$ — множество всех подмножеств $A$; $A \sqcup B$ — дизъюнктное объединение ($A \cap B = \emptyset$); $A \triangle B = (A \setminus B) \cup (B \setminus A)$ — симметрическая разность.
  \item $\iverson{P}$ — \textbf{скобка Айверсона}: $\iverson{P} = 1$, если высказывание $P$ истинно, и $\iverson{P} = 0$, если $P$ ложно; $0^0 = 1$, $0! = 1$.
  \item $d \mid n$ — целое $d$ делит $n$; $\gcd(a, b)$ и $\lcm(a, b)$ — наибольший общий делитель и наименьшее общее кратное.
  \item $\fall{x}{k} = x(x-1)\dots(x-k+1)$ — \textbf{убывающий факториал}; $\rise{x}{k} = x(x+1)\dots(x+k-1)$ — \textbf{возрастающий факториал} ($\fall{x}{0} = \rise{x}{0} = 1$).
  \item $S_n$ — симметрическая группа всех $n!$ перестановок множества $[n]$; $[x^n] F(x)$ — коэффициент при $x^n$ в формальном степенном ряде $F(x)$.
\end{itemize}
\end{summarybox}

{\small\tableofcontents}
\vspace{4pt}
\hrule
\vspace{6pt}

% ==============================================================================
\section{Уровень 1: База и модульная арифметика (CF 1300--1700)}
% ==============================================================================

\subsection{Модульная арифметика и базовая комбинаторика}

\textbf{Кольцо вычетов.} Для целого модуля $m \ge 2$ числа $a, b \in \Z$ сравнимы по модулю $m$ ($a \equiv b \pmod m$), если $m \mid (a - b)$. Канонический остаток обозначается $(a \bmod m) \in \{0, 1, \dots, m - 1\}$. Множество классов вычетов образует коммутативное кольцо $\Z_m$; при простом $p$ кольцо $\Z_p$ является конечным полем $\F_p$.

\textbf{Бинарное возведение в степень (\texttt{binpow}).} Вычисление $a^n \bmod m$ ($a \in \Z$, $n \in \N_0$) за $O(\log n)$ умножений по рекуррентному соотношению:
\[
  a^0 = 1, \qquad
  a^n \equiv
  \begin{cases}
    (a^{n/2})^2 \pmod m, & n \equiv 0 \pmod 2, \\
    a \cdot a^{n-1} \pmod m, & n \equiv 1 \pmod 2.
  \end{cases}
\]

\textbf{Обратный элемент по модулю.} Для $a \in \Z$ и $m \ge 2$ обратным по модулю $m$ называется такой вычет $a^{-1} \in \{1, \dots, m - 1\}$, что $a \cdot a^{-1} \equiv 1 \pmod m$. Обратный элемент существует и единствен тогда и только тогда, когда $\gcd(a, m) = 1$:
\begin{itemize}
  \item \textbf{Простой модуль $p$ ($p \nmid a$):} по малой теореме Ферма $a^{p-1} \equiv 1 \pmod p$, откуда $a^{-1} \equiv a^{p-2} \pmod p$ за $O(\log p)$.
  \item \textbf{Произвольный модуль $m$ ($\gcd(a, m) = 1$):} по теореме Эйлера $a^{-1} \equiv a^{\varphi(m)-1} \pmod m$, где $\varphi(m) = |\{1 \le x \le m : \gcd(x, m) = 1\}|$ — функция Эйлера, либо за $O(\log m)$ через расширенный алгоритм Евклида: из $ax + my = 1$ следует $a^{-1} \equiv x \pmod m$.
  \item \textbf{Линейный предподсчёт всех обратных $1^{-1}, \dots, N^{-1} \bmod p$} ($p$ простое, $N < p$) за $O(N)$:
  \[
    \inv[1] = 1, \qquad \inv[i] = \bigl(p - \lfloor p / i \rfloor \cdot \inv[p \bmod i] \bmod p\bigr) \bmod p \quad (i = 2, \dots, N).
  \]
\end{itemize}

\textbf{Линейный предподсчёт \texttt{fact[]} и \texttt{inv\_fact[]}.} Для простого модуля $p$ и $N < p$ таблицы факториалов $\mathrm{fact}[i] = i! \bmod p$ и обратных факториалов $\mathrm{inv\_fact}[i] = (i!)^{-1} \bmod p$ вычисляются за $O(N + \log p)$:
\[
  \mathrm{fact}[0] = 1,\quad \mathrm{fact}[i] = \mathrm{fact}[i-1] \cdot i \bmod p \;\;(1 \le i \le N);
\]
\[
  \mathrm{inv\_fact}[N] = (\mathrm{fact}[N])^{p-2} \bmod p,\quad
  \mathrm{inv\_fact}[i-1] = \mathrm{inv\_fact}[i] \cdot i \bmod p \;\;(i = N, N-1, \dots, 1).
\]
При известной таблице обратный элемент числа $i \in [1, N]$ равен $i^{-1} \equiv \mathrm{fact}[i-1] \cdot \mathrm{inv\_fact}[i] \pmod p$ за $O(1)$.

\textbf{Базовые комбинаторные конфигурации} (конечное множество $U$ из $n = |U|$ различных элементов, $k \in \N_0$):
\begin{itemize}
  \item \textbf{Правила суммы и произведения:} $|A \sqcup B| = |A| + |B|$; $|A \times B| = |A| \cdot |B|$.
  \item \textbf{Перестановки ($P_n$):} упорядоченные кортежи длины $n$ из всех элементов $U$ без повторений: $P_n = n! = \prod_{i=1}^n i$.
  \item \textbf{Перестановки с повторениями (мультиномиальный коэффициент):} число упорядоченных перестановок мультимножества мощности $n$ с кратностями типов $k_1, \dots, k_m$ ($\sum_{i=1}^m k_i = n$):
  \[
    \binom{n}{k_1, k_2, \dots, k_m} = \frac{n!}{k_1!\,k_2!\dots k_m!}.
  \]
  \item \textbf{Размещения без повторений ($A_n^k$) и с повторениями ($\bar{A}_n^k$):} упорядоченные кортежи длины $k$ из элементов $U$:
  \[
    A_n^k = \fall{n}{k} = \frac{n!}{(n - k)!} \quad (0 \le k \le n, \text{ иначе } 0), \qquad \bar{A}_n^k = n^k.
  \]
  \item \textbf{Сочетания без повторений ($C_n^k = \binom{n}{k}$):} неупорядоченные $k$-элементные подмножества множества $U$:
  \[
    \binom{n}{k} = \frac{n!}{k!\,(n - k)!} = \frac{\fall{n}{k}}{k!} \quad (0 \le k \le n), \qquad \binom{n}{k} = 0 \text{ при } k < 0 \text{ или } k > n.
  \]
  По простому модулю $p > n$: $\binom{n}{k} \equiv \mathrm{fact}[n] \cdot \mathrm{inv\_fact}[k] \cdot \mathrm{inv\_fact}[n-k] \pmod p$ за $O(1)$.
\end{itemize}

\textbf{Бином Ньютона и основные тождества.} Для элементов $a, b$ коммутативного кольца и $n \in \N_0$:
\[
  (a + b)^n = \sum_{k=0}^{n} \binom{n}{k} a^{n-k} b^k.
\]
Базовые свойства биномиальных коэффициентов:
\begin{itemize}
  \item \textbf{Симметрия и суммы строки:} $\binom{n}{k} = \binom{n}{n-k}$; $\sum_{k=0}^n \binom{n}{k} = 2^n$; $\sum_{k=0}^n (-1)^k \binom{n}{k} = \iverson{n = 0}$; $\sum_{k=0}^n k\binom{n}{k} = n 2^{n-1}$.
  \item \textbf{Поглощение и перенос множителя:} $k\binom{n}{k} = n\binom{n-1}{k-1}$ ($k \ge 1$); $\binom{n}{m}\binom{m}{k} = \binom{n}{k}\binom{n-k}{m-k}$.
  \item \textbf{Свёртка Вандермонда:} $\sum_{j=0}^{k} \binom{n}{j}\binom{m}{k-j} = \binom{n+m}{k}$; в частности, $\sum_{j=0}^n \binom{n}{j}^2 = \binom{2n}{n}$.
\end{itemize}

\textbf{Треугольник Паскаля.} Таблица биномиальных коэффициентов строится за $O(N^2)$ по \textbf{любому} (в том числе составному) модулю $M$ без операции деления:
\[
  \binom{n}{0} = \binom{n}{n} = 1, \qquad \binom{n}{k} = \binom{n-1}{k-1} + \binom{n-1}{k} \quad (1 \le k \le n - 1).
\]

\textbf{Тождество хоккейной клюшки (Hockey-Stick Identity).} Для любых $n, r, m \in \N_0$:
\[
  \sum_{i=r}^{n} \binom{i}{r} = \binom{n+1}{r+1} \quad (\text{при } n < r \text{ сумма равна } 0), \qquad
  \sum_{j=0}^{m} \binom{r+j}{j} = \binom{r+m+1}{m}.
\]

% ------------------------------------------------------------------------------
\subsection{Шары и перегородки (Stars and Bars)}

\textbf{Базовая модель (композиции и сочетания с повторениями).} Число способов разложить $n \in \N_0$ неразличимых шаров по $k \ge 1$ различимым ящикам совпадает с числом целочисленных векторов $(x_1, \dots, x_k) \in \Z^k$, удовлетворяющих уравнению $x_1 + x_2 + \dots + x_k = n$:
\begin{itemize}
  \item \textbf{Неотрицательные решения ($x_i \ge 0$, слабые $k$-композиции / сочетания с повторениями $\bar{C}_k^n$):}
  \[
    |\{x \in \Z^k : x_1 + \dots + x_k = n,\; x_i \ge 0\}| = \binom{n + k - 1}{k - 1} = \binom{n + k - 1}{n}.
  \]
  \item \textbf{Положительные решения ($x_i \ge 1$, строгие $k$-композиции):}
  \[
    |\{x \in \Z^k : x_1 + \dots + x_k = n,\; x_i \ge 1\}| = \binom{n - 1}{k - 1} \quad (n \ge k, \text{ иначе } 0).
  \]
\end{itemize}

\textbf{Сдвиг при ограничениях снизу ($x_i \ge a_i$).} Для произвольных $a_1, \dots, a_k \in \Z$ замена переменных $y_i = x_i - a_i \ge 0$ сводит уравнение к $y_1 + \dots + y_k = n'$, где $n' = n - \sum_{i=1}^k a_i$. Число решений равно $\binom{n' + k - 1}{k - 1}$ при $n' \ge 0$ и $0$ при $n' < 0$.

\textbf{Переход от неравенств к равенствам.} Введение фиктивной остаточной переменной $x_{k+1} = n - \sum_{i=1}^k x_i \ge 0$ устанавливает биекцию между решениями неравенства и уравнения от $k + 1$ переменных:
\[
  |\{x \in \Z^k : x_1 + \dots + x_k \le n,\; x_i \ge 0\}| = \binom{n + k}{k}.
\]
Строгое неравенство $x_1 + \dots + x_k < n$ в целых числах эквивалентно $x_1 + \dots + x_k \le n - 1$.

\textbf{Двусторонние ограничения через ФВИ ($a_i \le x_i \le b_i$).} После сдвига $x_i \leftarrow x_i - a_i$ задача имеет вид $x_1 + \dots + x_k = n$, $0 \le x_i \le c_i$ ($c_i = b_i - a_i \ge 0$). Пусть $A_i = \{x : x_i \ge c_i + 1\}$ — множество решений с нарушенной $i$-й верхней границей. По формуле включений-исключений число допустимых решений равно:
\[
  \sum_{S \subseteq [k]} (-1)^{|S|} \binom{n - \sum_{i \in S}(c_i + 1) + k - 1}{k - 1},
\]
где слагаемое с верхним индексом меньше $k - 1$ равно $0$. При одинаковых границах $0 \le x_i \le c$ формула свёртывается за $O\bigl(\min(k, \lfloor n/(c+1)\rfloor)\bigr)$:
\[
  \sum_{j=0}^{\min(k,\, \lfloor n / (c + 1) \rfloor)} (-1)^j \binom{k}{j} \binom{n - j(c + 1) + k - 1}{k - 1}.
\]

% ------------------------------------------------------------------------------
\subsection{Разложение Кантора и код Лемера}

\textbf{Код Лемера (вектор инверсий).} Для перестановки $\pi = (\pi_1, \dots, \pi_n) \in S_n$ её \textbf{кодом Лемера} называется целочисленный вектор $L(\pi) = (c_1, c_2, \dots, c_n)$, где $c_i$ — число элементов правее $i$, меньших $\pi_i$:
\[
  c_i = |\{j \in \{i + 1, \dots, n\} : \pi_j < \pi_i\}|, \qquad 0 \le c_i \le n - i \quad (c_n = 0).
\]
Отображение $L : S_n \to [0, n-1] \times [0, n-2] \times \dots \times \{0\}$ является биекцией; сумма $\sum_{i=1}^n c_i = \inv(\pi)$ равна числу инверсий перестановки $\pi$.

\textbf{Факториальная система счисления (разложение Кантора).} Любое целое $X \in [0, n! - 1]$ единственным образом записывается в виде:
\[
  X = \sum_{i=1}^{n} c_i \cdot (n - i)!, \qquad 0 \le c_i \le n - i \quad (c_n = 0).
\]
Последовательное извлечение цифр с младших разрядов: для $k = 1, 2, \dots, n$ полагаем $c_{n + 1 - k} = X \bmod k$, $X \leftarrow \lfloor X / k \rfloor$.

\textbf{Лексикографическая биекция.} Функция $\rk(\pi) = \sum_{i=1}^n c_i (n - i)! \in [0, n! - 1]$ задаёт $0$-индексированный лексикографический номер перестановки $\pi \in S_n$: $\pi <_{\mathrm{lex}} \sigma \iff \rk(\pi) < \rk(\sigma)$.

\textbf{Перевод $\pi \leftrightarrow L(\pi)$ за $O(n \log n)$ через дерево Фенвика (BIT).} Поддерживается $1$-индексированное дерево Фенвика $T[1 \dots n]$ индикаторов ещё не использованных элементов (в начале все $n$ позиций равны $1$):
\begin{itemize}
  \item \textbf{Прямой ход $\pi \mapsto (c_1, \dots, c_n)$:} для $i = 1, \dots, n$ величина $c_i$ равна сумме на префиксе $[1, \pi_i - 1]$ в $T$, после чего выполняется обновление $\mathrm{add}(\pi_i, -1)$. Время: $O(n \log n)$.
  \item \textbf{Обратный ход $(c_1, \dots, c_n) \mapsto \pi$:} для $i = 1, \dots, n$ элемент $\pi_i$ является $(c_i + 1)$-м по величине среди доступных, то есть наименьшим индексом $v \in [1, n]$ с префиксной суммой $\ge c_i + 1$. Поиск $v$ выполняется бинарным спуском по дереву Фенвика за $O(\log n)$: полагаем $\mathrm{pos} = 0$, $\mathrm{rem} = c_i + 1$; для $b = 2^{\lfloor \log_2 n \rfloor}, \dots, 2, 1$, если $\mathrm{pos} + b \le n$ и $T[\mathrm{pos} + b] < \mathrm{rem}$, увеличиваем $\mathrm{pos} \mathrel{+}= b$ и уменьшаем $\mathrm{rem} \mathrel{-}= T[\mathrm{pos}]$. Тогда $\pi_i = \mathrm{pos} + 1$, после чего вызывается $\mathrm{add}(\pi_i, -1)$. Суммарное время: $O(n \log n)$.
\end{itemize}

% ------------------------------------------------------------------------------
\subsection{Базовый тервер}

\textbf{Дискретное вероятностное пространство.} Тройка $(\Omega, \mathcal{F}, P)$, где $\Omega$ — не более чем счётное множество элементарных исходов, $\mathcal{F} = 2^\Omega$ — множество событий, $P : \mathcal{F} \to [0, 1]$ — вероятностная мера: $P(\Omega) = 1$, $P(\emptyset) = 0$ и $P\bigl(\bigsqcup_i A_i\bigr) = \sum_i P(A_i)$ для попарно несовместных событий ($A_i \cap A_j = \emptyset$ при $i \ne j$). Для любых $A, B \subseteq \Omega$:
\[
  P(\bar{A}) = 1 - P(A) \;\; (\bar{A} = \Omega \setminus A), \qquad P(A \cup B) = P(A) + P(B) - P(A \cap B).
\]

\textbf{Классическая вероятность.} Если пространство $\Omega$ конечно ($|\Omega| < \infty$) и все элементарные исходы равновероятны ($P(\{\omega\}) = 1 / |\Omega|$), то для любого события $A \subseteq \Omega$:
\[
  P(A) = \frac{|A|}{|\Omega|}.
\]

\textbf{Условная вероятность и теорема умножения.} При $P(B) > 0$ \textbf{условной вероятностью} события $A$ при наступлении $B$ называется величина $P(A \mid B) = \frac{P(A \cap B)}{P(B)}$. Цепное правило для событий с $P(A_1 \cap \dots \cap A_{k-1}) > 0$:
\[
  P(A_1 \cap A_2 \cap \dots \cap A_k) = P(A_1)\, P(A_2 \mid A_1)\, P(A_3 \mid A_1 \cap A_2) \dots P(A_k \mid A_1 \cap \dots \cap A_{k-1}).
\]

\textbf{Независимость событий.}
\begin{itemize}
  \item Два события $A$ и $B$ \textbf{независимы}, если $P(A \cap B) = P(A)P(B)$ (при $P(B) > 0$ равносильно $P(A \mid B) = P(A)$).
  \item События $A_1, \dots, A_n$ \textbf{независимы в совокупности}, если для любого непустого подмножества индексов $I \subseteq [n]$ выполнено $P\bigl(\bigcap_{i \in I} A_i\bigr) = \prod_{i \in I} P(A_i)$. Из попарной независимости независимость в совокупности не следует.
\end{itemize}

\textbf{Формула полной вероятности и формула Байеса.} Пусть $H_1, \dots, H_k$ — \textbf{полная группа гипотез}: $\Omega = \bigsqcup_{i=1}^k H_i$ и $P(H_i) > 0$ для всех $i$. Тогда для любого события $A$:
\[
  P(A) = \sum_{i=1}^{k} P(H_i)\, P(A \mid H_i), \qquad
  P(H_j \mid A) = \frac{P(H_j)\, P(A \mid H_j)}{\sum_{i=1}^{k} P(H_i)\, P(A \mid H_i)} \quad \bigl(P(A) > 0\bigr).
\]

\textbf{Подсчёт вероятностей по простому модулю.} Пусть ответ (вероятность или матожидание) равен несократимой рациональной дроби $P / Q$ ($P \in \Z$, $Q \in \N$, $\gcd(P, Q) = 1$), и фиксирован простой модуль $p$ (например, $10^9 + 7$ или $998244353$), не делящий $Q$ ($p \nmid Q$). Дроби сопоставляется вычет:
\[
  \frac{P}{Q} \bmod p \;:=\; (P \bmod p) \cdot Q^{p-2} \bmod p \in \F_p.
\]
Отображение локального кольца $\Z_{(p)} = \{a / b \in \mathbb{Q} : p \nmid b\} \to \F_p$ является гомоморфизмом колец, поэтому сложение, вычитание, умножение и деление рациональных величин можно выполнять непосредственно по модулю $p$ на всех промежуточных этапах при условии, что ни один промежуточный знаменатель не кратен $p$.

% ==============================================================================
\section{Уровень 2: Рабочая классика контестов (CF 1700--2000)}
% ==============================================================================

\subsection{Линейность математического ожидания}

\textbf{Случайная величина и математическое ожидание.} \textbf{Дискретная случайная величина (СВ)} — функция $X : \Omega \to \R$ с не более чем счётным множеством значений $\operatorname{Im}(X)$. При условии абсолютной сходимости $\sum_{x} |x| P(X = x) < \infty$ её \textbf{математическим ожиданием} называется среднее взвешенное значение:
\[
  \E[X] = \sum_{\omega \in \Omega} X(\omega) P(\{\omega\}) = \sum_{x \in \operatorname{Im}(X)} x \, P(X = x), \qquad
  \E[g(X)] = \sum_{x \in \operatorname{Im}(X)} g(x) \, P(X = x).
\]

\textbf{Линейность матожидания.} Для \textbf{произвольных} (в том числе как угодно зависимых) случайных величин $X_1, \dots, X_n$ с конечными матожиданиями и констант $c_1, \dots, c_n \in \R$:
\[
  \E\!\left[\sum_{i=1}^{n} c_i X_i\right] = \sum_{i=1}^{n} c_i \E[X_i].
\]
Если СВ $X$ и $Y$ \textbf{независимы} ($P(X = x, Y = y) = P(X = x)P(Y = y)$ для всех $x, y$), то дополнительно $\E[XY] = \E[X]\E[Y]$. \textbf{Дисперсия} $\Var(X) = \E[(X - \E[X])^2] = \E[X^2] - (\E[X])^2$ аддитивна для попарно независимых СВ.

\textbf{Метод индикаторов.} \textbf{Индикатор} события $A \subseteq \Omega$ — случайная величина $\Ind_A(\omega) = \iverson{\omega \in A} \in \{0, 1\}$. Основные свойства:
\[
  \E[\Ind_A] = P(A), \qquad \Ind_A^k = \Ind_A \;\;(k \ge 1), \qquad \Ind_{A \cap B} = \Ind_A \cdot \Ind_B.
\]
Если величина $X$ считает количество наступивших событий из набора $A_1, \dots, A_m$, то $X = \sum_{i=1}^m \Ind_{A_i}$ и
\[
  \E[X] = \sum_{i=1}^{m} P(A_i), \qquad
  \E[X^2] = \sum_{i=1}^{m} P(A_i) + 2 \sum_{1 \le i < j \le m} P(A_i \cap A_j).
\]

\textbf{Хвостовая формула (Tail Sum Formula).} Для любой случайной величины $X$, принимающей значения в $\N_0 = \{0, 1, 2, \dots\}$:
\[
  \E[X] = \sum_{k=1}^{\infty} P(X \ge k) = \sum_{k=0}^{\infty} P(X > k), \qquad
  \E[X^2] = \sum_{k=1}^{\infty} (2k - 1) P(X \ge k).
\]

\textbf{Геометрическое распределение $T \sim \Geom(p)$.} В последовательности независимых испытаний Бернулли с вероятностью успеха $p \in (0, 1]$ случайная величина $T \in \N$ — номер испытания до первого успеха включительно:
\[
  P(T = k) = (1 - p)^{k-1} p \;\;(k \ge 1), \qquad P(T \ge k) = (1 - p)^{k-1}, \qquad
  \E[T] = \frac{1}{p}, \qquad \Var(T) = \frac{1 - p}{p^2}.
\]
Свойство отсутствия памяти: $P(T > m + k \mid T > m) = P(T > k)$ для всех $m, k \in \N_0$.

\textbf{Задача о собирателе купонов (Coupon Collector's Problem).} На каждом шаге независимо и равновероятно выбирается один из $n$ типов купонов. Пусть $T$ — число шагов до появления всех $n$ типов. Представим $T = \sum_{i=0}^{n-1} T_i$, где $T_i \sim \Geom\bigl(\frac{n - i}{n}\bigr)$ — время ожидания $(i + 1)$-го нового типа после сбора $i$ различных типов (величины $T_i$ независимы):
\[
  \E[T] = \sum_{i=0}^{n-1} \frac{n}{n - i} = n H_n = n \sum_{k=1}^{n} \frac{1}{k} = n \ln n + \gamma n + O(1),
\]
где $H_n$ — $n$-е гармоническое число, $\gamma \approx 0{,}57721566$ — постоянная Эйлера—Маскерони. Ожидаемое число шагов до сбора первых $m \le n$ различных типов равно $n(H_n - H_{n-m})$.

% ------------------------------------------------------------------------------
\subsection{Формула включений-исключений (ФВИ / PIE)}

\textbf{Общая формулировка.} Пусть $U$ — конечный универсум и $A_1, \dots, A_k \subseteq U$ — семейство свойств (подмножеств). Для мощности объединения и мощности пересечения дополнений ($\bar{A}_i = U \setminus A_i$, при соглашении $\bigcap_{i \in \emptyset} A_i = U$):
\[
  \left|\bigcup_{i=1}^{k} A_i\right| = \sum_{\emptyset \ne S \subseteq [k]} (-1)^{|S| - 1} \left|\bigcap_{i \in S} A_i\right|, \qquad
  \left|\bigcap_{i=1}^{k} \bar{A}_i\right| = \sum_{S \subseteq [k]} (-1)^{|S|} \left|\bigcap_{i \in S} A_i\right|.
\]
В вероятностной форме: $P\bigl(\bigcup_{i=1}^k A_i\bigr) = \sum_{\emptyset \ne S \subseteq [k]} (-1)^{|S|-1} P\bigl(\bigcap_{i \in S} A_i\bigr)$. Если размер пересечения $\left|\bigcap_{i \in S} A_i\right| = f(|S|)$ зависит только от мощности $j = |S|$, то $\left|\bigcap_{i=1}^k \bar{A}_i\right| = \sum_{j=0}^k (-1)^j \binom{k}{j} f(j)$.

\textbf{Перебор битовых масок за $O(2^k)$.} При малом $k$ ($k \le 20\text{--}25$) подмножество $S \subseteq [k]$ кодируется маской $m \in [0, 2^k - 1]$, а знак слагаемого равен $(-1)^{\popcnt(m)}$. При пересчёте пересечения через отщепление младшего единичного бита $b = m \mathbin{\&} (-m)$ или рекурсивным обходом дерева подмножеств суммарное время составляет $O(2^k)$ вместо $O(k \cdot 2^k)$.

\textbf{Беспорядки (субфакториал $!n$).} Перестановка $\pi \in S_n$ называется \textbf{беспорядком} (derangement), если она не имеет неподвижных точек ($\pi_i \ne i$ для всех $i \in [n]$). Число беспорядков обозначается $!n$ (или $D_n$):
\[
  !n = n! \sum_{j=0}^{n} \frac{(-1)^j}{j!}, \qquad !0 = 1, \quad !1 = 0, \quad !2 = 1.
\]
Линейные рекурренты для вычисления $!0, \dots, !N \bmod M$ за $O(N)$ по произвольному модулю $M$:
\[
  !n = (n - 1)\bigl(!(n - 1) + !(n - 2)\bigr) \;\; (n \ge 2), \qquad
  !n = n \cdot !(n - 1) + (-1)^n \;\; (n \ge 1).
\]
При $n \ge 1$ выполнено $!n = \bigl\lfloor \frac{n!}{e} + \frac{1}{2} \bigr\rfloor$ и $\lim_{n \to \infty} \frac{!n}{n!} = e^{-1}$. Число перестановок в $S_n$ с ровно $m$ неподвижными точками равно $\binom{n}{m} \cdot !(n - m)$.

\textbf{ФВИ внутри динамического программирования.} Если число условий $k$ велико ($k \sim n$), но вклад пересечения $\left|\bigcap_{i \in S} A_i\right|$ распадается в произведение локальных переходов вдоль объекта, явное хранение маски $S$ не требуется. В состоянии ДП фиксируется лишь вспомогательный параметр, определяющий вес пересечения, а знак $(-1)^{|S|}$ накапливается одним из двух эквивалентных способов:
\begin{enumerate}
  \item \textbf{Чётность в состоянии:} добавляется булев флаг $p = |S| \bmod 2 \in \{0, 1\}$ (число уже выбранных «нарушенных» условий по модулю $2$); при включении очередного условия в $S$ чётность инвертируется $p \leftarrow p \oplus 1$, а итоговый ответ равен $\mathrm{dp}_{\text{end}}[0] - \mathrm{dp}_{\text{end}}[1]$.
  \item \textbf{Свёртывание знака в коэффициентах переходов:} каждый выбор очередного нарушенного условия домножает вклад перехода на $(-1)$ (или блок длины $L$, склеивающий $L - 1$ нарушений подряд, входит с весом $(-1)^{L-1}$), что полностью убирает размерность $S$ из состояния.
\end{enumerate}

% ------------------------------------------------------------------------------
\subsection{Числа Каталана и принцип отражений (Андре)}

\textbf{Числа Каталана ($C_n$, $n \ge 0$).} Последовательность $(C_n)_{n \ge 0} = (1, 1, 2, 5, 14, 42, 132, \ldots)$:
\[
  C_n = \frac{1}{n + 1}\binom{2n}{n} = \binom{2n}{n} - \binom{2n}{n + 1} \quad (n \ge 0).
\]
Рекуррентные соотношения (квадратичная свёртка по любому модулю $M$ и линейный пересчёт):
\[
  C_0 = 1, \qquad C_n = \sum_{i=0}^{n-1} C_i C_{n - 1 - i} \;\;(n \ge 1), \qquad
  C_n = \frac{2(2n - 1)}{n + 1}\, C_{n - 1}.
\]

\textbf{Классические комбинаторные объекты, перечисляемые числом $C_n$:}
\begin{enumerate}
  \item \textbf{Правильные скобочные последовательности (ПСП)} длины $2n$ из $n$ открывающих и $n$ закрывающих скобок (на любом префиксе баланс $\ge 0$, полный баланс равен $0$).
  \item \textbf{Пути Дика на решётке:} монотонные пути из $(0, 0)$ в $(n, n)$ с шагами $(1, 0)$ и $(0, 1)$, не поднимающиеся выше диагонали $y = x$ ($y \le x$ во всех точках); эквивалентно пути из $(0, 0)$ в $(2n, 0)$ с шагами $(1, +1)$ и $(1, -1)$ в полуплоскости $y \ge 0$.
  \item \textbf{Плоские бинарные и корневые деревья:} упорядоченные бинарные деревья с $n$ внутренними вершинами (эквивалентно полные бинарные деревья с $n + 1$ листьями, либо произвольные плоские корневые деревья на $n + 1$ вершинах).
  \item \textbf{Триангуляции выпуклого $(n + 2)$-угольника:} разбиения выпуклого многоугольника с помеченными вершинами на $n$ треугольников с помощью $n - 1$ непересекающихся внутренних диагоналей.
  \item \textbf{Непересекающиеся хорды и стековые перестановки:} способы соединить $2n$ точек на окружности $n$ попарно непересекающимися хордами; перестановки из $S_n$, сортируемые одним стеком (избегающие паттерна $231$).
\end{enumerate}

\textbf{Принцип отражений Андре (одна запрещённая прямая).} Число всех монотонных путей на решётке из $(x_0, y_0)$ в $(x_1, y_1)$ с шагами $(1, 0)$ и $(0, 1)$ равно $\binom{\Delta x + \Delta y}{\Delta x}$, где $\Delta x = x_1 - x_0 \ge 0$, $\Delta y = y_1 - y_0 \ge 0$. Пусть задана запрещённая прямая $L : y = x + c$ ($c \in \Z$), и обе концевые точки лежат строго ниже $L$: $y_0 < x_0 + c$, $y_1 < x_1 + c$. Отражение точки $(x, y)$ относительно $L$ переводит её в $(y - c, x + c)$ и устанавливает биекцию между путями из $(x_0, y_0)$ в $(x_1, y_1)$, имеющими хотя бы одну общую точку с $L$, и всеми путями из отражённого старта $(y_0 - c, x_0 + c)$ в $(x_1, y_1)$. Число путей, лежащих строго ниже $L$ (то есть удовлетворяющих $y \le x + c - 1$):
\[
  N_{< L} = \binom{\Delta x + \Delta y}{\Delta x} - \binom{\Delta x + \Delta y}{\Delta x + c - (y_0 - x_0)}.
\]
\emph{Следствие (обобщённая формула Каталана / задача о баллотировке):} число путей из $(0, 0)$ в $(n, m)$ ($n \ge m \ge 0$), не поднимающихся выше $y = x$ ($c = 1$), равно $\binom{n + m}{n} - \binom{n + m}{n + 1} = \frac{n - m + 1}{n + 1}\binom{n + m}{m}$.

\textbf{Обобщение: пути строго между двумя параллельными прямыми.} Пусть требуется посчитать монотонные пути из $(x_0, y_0)$ в $(x_1, y_1)$, не заходящие на запрещённые прямые $L_1 : y = x + a$ и $L_2 : y = x - b$ ($a, b \in \Z$, $x_0 - b < y_0 < x_0 + a$, $x_1 - b < y_1 < x_1 + a$; для нестрогого коридора $x - b' \le y \le x + a'$ полагают $a = a' + 1$, $b = b' + 1$). Положим $D = a + b \ge 2$ — ширину полосы. Многократные отражения относительно $L_1$ и $L_2$ с учётом ФВИ по чередующимся касаниям $L_1, L_2, \dots$ дают точную формулу за $O\bigl(\frac{\Delta x + \Delta y}{D}\bigr)$ ненулевых слагаемых:
\[
  N_{L_2 < \text{путь} < L_1} = \sum_{k \in \Z} \binom{\Delta x + \Delta y}{\Delta x - kD} - \sum_{k \in \Z} \binom{\Delta x + \Delta y}{\Delta x + a - (y_0 - x_0) + kD}, \qquad D = a + b.
\]

% ------------------------------------------------------------------------------
\subsection{ДП по вероятностям и матожиданиям}

\textbf{Постановка задачи на пространстве состояний.} Задано конечное множество состояний $V$, начальное состояние $s \in V$ и множество терминальных (поглощающих) состояний $T \subseteq V$. Из каждого $u \in V \setminus T$ процесс переходит в $v \in V$ с вероятностью $p_{u,v} \ge 0$ ($\sum_{v \in V} p_{u,v} = 1$) и платой за шаг $c_{u,v}$ (для подсчёта числа шагов $c_{u,v} = 1$). Обозначим среднюю стоимость одного шага из $u$ через $\bar{c}_u = \sum_{v \in V} p_{u,v} c_{u,v}$.

\textbf{Прямое ДП (Forward DP) и ловушка условного матожидания.} Используется на ациклическом графе переходов (в порядке топологической сортировки) для подсчёта вероятности $P[u]$ посетить состояние $u$:
\[
  P[s] = 1, \qquad P[v] = \sum_{u : u \to v} P[u] \cdot p_{u,v}.
\]
Для матожиданий наивный пересчёт $\sum_{u \to v} p_{u,v}(E[u] + c_{u,v})$ некорректен, так как переходы в $v$ условны по факту попадания в $v$. При прямом обходе поддерживается пара $(P[u], S[u])$, где $S[u] = P[u] \cdot \E[\text{стоимость до } u \mid \text{попали в } u]$ — взвешенная сумма стоимостей путей:
\[
  S[s] = 0, \qquad S[v] = \sum_{u : u \to v} \bigl(S[u] + P[u] \cdot c_{u,v}\bigr) p_{u,v}, \qquad
  \E[\text{стоимость до } v \mid \text{попали в } v] = \frac{S[v]}{P[v]}.
\]

\textbf{Обратное ДП (Backward DP — стандарт для матожиданий).} Определим $E[u]$ как ожидаемую суммарную стоимость пути от старта в состоянии $u$ до первого попадания в $T$ (при условии достижимости $T$ с вероятностью $1$), а $f[u]$ — вероятность завершить процесс в целевом подмножестве $T_{\text{win}} \subseteq T$:
\[
  E[t] = 0 \;\; (t \in T), \qquad E[u] = \bar{c}_u + \sum_{v \in V} p_{u,v} E[v] \;\; (u \notin T);
\]
\[
  f[t] = \iverson{t \in T_{\text{win}}} \;\; (t \in T), \qquad f[u] = \sum_{v \in V} p_{u,v} f[v] \;\; (u \notin T).
\]
На DAG пересчёт идёт в обратном топологическом порядке; ответом для старта служит $E[s]$ (соответственно $f[s]$) без нормирующих делений на $P[u]$.

\textbf{Избавление от петель ($p_{u,u} > 0$).} Если все циклы в графе переходов имеют длину $1$ (переходы вида $u \to u$ с $0 \le p_{u,u} < 1$), слагаемое с $E[u]$ переносится в левую часть уравнения:
\[
  E[u] = \frac{\bar{c}_u + \sum_{v \ne u} p_{u,v} E[v]}{1 - p_{u,u}}, \qquad
  f[u] = \frac{\sum_{v \ne u} p_{u,v} f[v]}{1 - p_{u,u}}.
\]
В частности, при $c_{u,v} = 1$ числитель равен $1 + \sum_{v \ne u} p_{u,v} E[v]$.

% ==============================================================================
\section{Уровень 3: Комбинаторные числа, графы и СЛАУ (CF 2000--2300)}
% ==============================================================================

\subsection{Числа Стирлинга и числа Белла}

\textbf{Числа Стирлинга 1-го рода.} \textbf{Беззнаковое число Стирлинга 1-го рода} $\stirlingfirst{n}{k}$ (также $c(n, k)$, $0 \le k \le n$) равно числу перестановок в $S_n$, разлагающихся в произведение ровно $k$ независимых циклов:
\[
  \stirlingfirst{0}{0} = 1, \qquad \stirlingfirst{n}{0} = \stirlingfirst{0}{k} = 0 \;\;(n, k \ge 1), \qquad
  \stirlingfirst{n}{k} = (n - 1)\stirlingfirst{n - 1}{k} + \stirlingfirst{n - 1}{k - 1}.
\]
\textbf{Знаковое число Стирлинга 1-го рода} определяется как $s(n, k) = (-1)^{n - k}\stirlingfirst{n}{k}$. Алгебраический смысл — разложение возрастающих и убывающих факториалов по степеням $x$:
\[
  \rise{x}{n} = \sum_{k=0}^{n} \stirlingfirst{n}{k} x^k, \qquad
  \fall{x}{n} = \sum_{k=0}^{n} (-1)^{n - k} \stirlingfirst{n}{k} x^k = \sum_{k=0}^{n} s(n, k) x^k, \qquad
  \sum_{k=0}^{n} \stirlingfirst{n}{k} = n!.
\]

\textbf{Числа Стирлинга 2-го рода.} Величина $\stirlingsecond{n}{k}$ (также $S(n, k)$, $0 \le k \le n$) равна числу разбиений $n$-элементного множества на $k$ непустых попарно непересекающихся \textbf{неупорядоченных} блоков (число сюръекций $[n] \to [k]$ равно $k! \stirlingsecond{n}{k}$):
\[
  \stirlingsecond{0}{0} = 1, \qquad \stirlingsecond{n}{0} = \stirlingsecond{0}{k} = 0 \;\;(n, k \ge 1), \qquad
  \stirlingsecond{n}{k} = k\stirlingsecond{n - 1}{k} + \stirlingsecond{n - 1}{k - 1}.
\]
Явная формула через ФВИ (позволяет вычислить одно значение за $O(k \log n)$ или всю строку $\stirlingsecond{n}{0 \dots n}$ одной свёрткой полиномов за $O(n \log n)$):
\[
  \stirlingsecond{n}{k} = \frac{1}{k!} \sum_{j=0}^{k} (-1)^{k - j} \binom{k}{j} j^n = \sum_{j=0}^{k} \frac{(-1)^{k - j}}{(k - j)!} \cdot \frac{j^n}{j!}.
\]

\textbf{Числа Белла ($B_n$).} Общее число разбиений $n$-элементного множества на произвольное число непустых неупорядоченных блоков ($B_0 = 1, B_1 = 1, B_2 = 2, B_3 = 5, B_4 = 15, \dots$; для оценки перебора: $B_{10} = 115\,975$, $B_{12} = 4\,213\,597$):
\[
  B_n = \sum_{k=0}^{n} \stirlingsecond{n}{k}, \qquad
  B_{n + 1} = \sum_{k=0}^{n} \binom{n}{k} B_k, \qquad
  \sum_{n=0}^{\infty} B_n \frac{x^n}{n!} = \exp(e^x - 1).
\]
Сравнение Ташарда для простого модуля $p$: $B_{n + p} \equiv B_n + B_{n + 1} \pmod p$.

\textbf{Связь обычных степеней и убывающих факториалов; суммы вида $\sum_i f(i)^k$.} Для любого $x \in \R$ и целого $k \ge 0$ справедливо разложение (матрицы $\bigl[s(n, k)\bigr]$ и $\bigl[\stirlingsecond{n}{k}\bigr]$ взаимно обратны):
\[
  x^k = \sum_{j=0}^{k} \stirlingsecond{k}{j} \fall{x}{j} = \sum_{j=0}^{k} \stirlingsecond{k}{j} j! \binom{x}{j}.
\]
Применения при малом показателе $k$ и большом числе слагаемых:
\begin{itemize}
  \item \textbf{Сумма $k$-х степеней первых натуральных чисел} за $O(k^2)$ (через тождество хоккейной клюшки):
  \[
    \sum_{i=0}^{n} i^k = \sum_{j=0}^{k} \stirlingsecond{k}{j} j! \binom{n + 1}{j + 1} = \sum_{j=0}^{k} \stirlingsecond{k}{j} \frac{\fall{(n + 1)}{j + 1}}{j + 1}.
  \]
  \item \textbf{Биномиально-степенная сумма} за $O(k)$ (после предподсчёта строки Стирлинга):
  \[
    \sum_{i=0}^{n} \binom{n}{i} i^k = \sum_{j=0}^{k} \stirlingsecond{k}{j} \fall{n}{j} 2^{n - j}.
  \]
  \item \textbf{ДП для сумм $\sum_{\omega \in \mathcal{X}} f(\omega)^k$ (на деревьях, путях, подмножествах):} вместо степеней хранятся биномиальные моменты $D_j = \sum_{\omega \in \mathcal{X}} \binom{f(\omega)}{j}$ для $0 \le j \le k$. Увеличение $f \mapsto f + 1$ пересчитывается по правилу Паскаля $D'_j = D_j + D_{j-1}$ за $O(k)$, слияние независимых слагаемых $f_1 + f_2$ — свёрткой Вандермонда $D_j = \sum_{t=0}^j D^{(1)}_t D^{(2)}_{j - t}$ за $O(k^2)$, а искомая сумма равна $\sum_{j=0}^k \stirlingsecond{k}{j} j! \, D_j$.
\end{itemize}

% ------------------------------------------------------------------------------
\subsection{Теорема Люка и китайская теорема об остатках (КТО / CRT)}

\textbf{Теорема Люка (простой модуль $p$).} Пусть $p$ — простое число, а $n, k \in \N_0$ записаны в $p$-ичной системе счисления: $n = \sum_{i=0}^s n_i p^i$, $k = \sum_{i=0}^s k_i p^i$, где $0 \le n_i, k_i < p$. Тогда:
\[
  \binom{n}{k} \equiv \prod_{i=0}^{s} \binom{n_i}{k_i} \pmod p.
\]
В частности, $\binom{n}{k} \equiv 0 \pmod p \iff \exists i : k_i > n_i$; при $p = 2$: $\binom{n}{k} \equiv 1 \pmod 2 \iff (k \mathbin{\&} n) = k$. Предподсчёт факториалов до $p - 1$ за $O(p)$ позволяет вычислять $\binom{n}{k} \bmod p$ за $O(\log_p n)$ при $n \le 10^{18}$.

\textbf{Китайская теорема об остатках (CRT).} Пусть модули $m_1, \dots, m_r \ge 2$ попарно взаимно просты ($\gcd(m_i, m_j) = 1$ при $i \ne j$) и $M = \prod_{i=1}^r m_i$. Система сравнений $x \equiv a_i \pmod{m_i}$ ($1 \le i \le r$) имеет единственное решение по модулю $M$:
\[
  x \equiv \sum_{i=1}^{r} a_i M_i y_i \pmod M, \qquad M_i = \frac{M}{m_i}, \qquad y_i = M_i^{-1} \bmod m_i.
\]
\emph{Обобщённая КТО:} система из двух сравнений $x \equiv a_1 \pmod{m_1}$, $x \equiv a_2 \pmod{m_2}$ совместна тогда и только тогда, когда $a_1 \equiv a_2 \pmod g$, где $g = \gcd(m_1, m_2)$; при совместности решение единственно по модулю $\lcm(m_1, m_2)$ и находится из соотношения Безу $\frac{m_1}{g} u + \frac{m_2}{g} v = 1$ как $x = a_1 + m_1 \cdot \bigl(\frac{a_2 - a_1}{g} u \bmod \frac{m_2}{g}\bigr)$.

\textbf{Вычисление $\binom{n}{k} \bmod M$ при $n \le 10^{18}$ и составном модуле $M = \prod_{i=1}^r p_i^{\alpha_i}$:}
\begin{enumerate}
  \item \textbf{Модуль свободен от квадратов ($\alpha_i = 1$ для всех $i$):} вычисляем $r_i = \binom{n}{k} \bmod p_i$ по теореме Люка и восстанавливаем ответ по модулю $M$ через КТО.
  \item \textbf{Примарные компоненты $p^\alpha$ ($\alpha \ge 2$, метод Грэнвилла / $p$-адическое разложение факториала):}
  \begin{itemize}
    \item По \textbf{формуле Лежандра} показатель вхождения простого $p$ в $m!$ равен $v_p(m!) = \sum_{j \ge 1} \lfloor m / p^j \rfloor$. Показатель $e = v_p\bigl(\binom{n}{k}\bigr) = v_p(n!) - v_p(k!) - v_p((n - k)!)$ равен числу переносов при сложении $k + (n - k)$ в базе $p$ (\textbf{теорема Куммера}). Если $e \ge \alpha$, то $\binom{n}{k} \equiv 0 \pmod{p^\alpha}$.
    \item При $e < \alpha$ положим $(m!)_p = \prod_{1 \le j \le m,\, p \nmid j} j \bmod p^\alpha$ и представим $m! = p^{v_p(m!)} f_p(m)$, где $\gcd(f_p(m), p) = 1$. По периодичности взаимно простых вычетов с периодом $p^\alpha$:
    \[
      (m!)_p \equiv U^{\lfloor m / p^\alpha \rfloor} \cdot \bigl((m \bmod p^\alpha)!\bigr)_p \pmod{p^\alpha}, \qquad
      U = (p^\alpha!)_p \equiv \begin{cases} 1 \pmod{p^\alpha}, & p = 2,\; \alpha \ge 3, \\ -1 \pmod{p^\alpha}, & \text{иначе.} \end{cases}
    \]
    Тогда взаимно простая с $p$ часть факториала равна $f_p(m) \equiv \prod_{j \ge 0} \bigl(\lfloor m / p^j \rfloor !\bigr)_p \pmod{p^\alpha}$ и вычисляется за $O(\log_p m)$ при предподсчёте префиксных произведений длины $p^\alpha$ за $O(p^\alpha)$.
    \item Обращая $f_p(k)$ и $f_p(n - k)$ расширенным алгоритмом Евклида по модулю $p^\alpha$, получаем:
    \[
      \binom{n}{k} \equiv p^e \cdot f_p(n) \cdot \bigl(f_p(k)\bigr)^{-1} \cdot \bigl(f_p(n - k)\bigr)^{-1} \pmod{p^\alpha},
    \]
    после чего остатки по всем $p_i^{\alpha_i}$ объединяются через КТО.
  \end{itemize}
\end{enumerate}

% ------------------------------------------------------------------------------
\subsection{Цепи Маркова и метод Гаусса}

\textbf{Конечная однородная цепь Маркова.} Случайный процесс $(X_t)_{t \ge 0}$ на конечном множестве состояний $V = \{1, \dots, N\}$, обладающий \textbf{марковским свойством}:
\[
  P(X_{t + 1} = j \mid X_t = i, X_{t - 1} = i_{t - 1}, \dots, X_0 = i_0) = P(X_{t + 1} = j \mid X_t = i) = p_{i,j}.
\]
Матрица переходов $P = (p_{i,j})_{N \times N}$ стохастическая по строкам ($p_{i,j} \ge 0$, $\sum_{j=1}^N p_{i,j} = 1$). Если вектор-строка $\pi^{(0)}$ задаёт начальное распределение, то через $t$ шагов $\pi^{(t)} = \pi^{(0)} P^t$ (вычисляется за $O(N^3 \log t)$ матричным возведением в степень). \textbf{Стационарное распределение} $\pi$ удовлетворяет системе $\pi P = \pi$, $\sum_{i=1}^N \pi_i = 1$, $\pi_i \ge 0$; для неприводимой апериодической цепи оно единственно, и $\pi_i = 1 / \E_i[\tau_i^+]$, где $\tau_i^+$ — время первого возвращения в $i$.

\textbf{Циклическое ДП по матожиданиям через СЛАУ за $O(N^3)$.} Пусть $T \subseteq V$ — поглощающие состояния ($E[t] = 0$ при $t \in T$), а $Q = V \setminus T$ ($|Q| = m$) — невозвратные состояния, из каждого из которых $T$ достижимо с положительной вероятностью, причём подграф на $Q$ содержит циклы. Уравнения обратного ДП образуют СЛАУ размера $m \times m$:
\[
  E[u] - \sum_{v \in Q} p_{u,v} E[v] = \bar{c}_u \quad (u \in Q) \quad \Longleftrightarrow \quad (I - P_{QQ}) \mathbf{E} = \bar{\mathbf{c}}.
\]
Спектральный радиус подматрицы переходов $P_{QQ}$ строго меньше $1$, поэтому матрица $(I - P_{QQ})$ невырождена (фундаментальная матрица $F = (I - P_{QQ})^{-1} = \sum_{k=0}^\infty P_{QQ}^k$, где $F_{u,v}$ — ожидаемое число посещений состояния $v$ до поглощения при старте из $u$). Система решается методом Гаусса за $O(m^3)$ (с выбором главного элемента в вещественной арифметике или над полем $\F_p$).

\textbf{Оптимизации метода Гаусса для разреженных структур:}
\begin{itemize}
  \item \textbf{Ленточные матрицы за $O(N W^2)$:} если все ненулевые коэффициенты системы сосредоточены в полосе $|i - j| \le W$ (например, блуждание на таблице $H \times W$, занумерованной по строкам длинной стороны, или одномерные скачки длины $\le W$), в прямом ходе Гаусса исключаются только $W$ строк ниже диагонали с пересчётом не более $2W$ столбцов (с учётом частичного пивотирования внутри ленты), что даёт время $O(N W^2)$ и память $O(N W)$. При $W = 1$ (трёхдиагональная система, 1D-блуждание) применяется \textbf{метод прогонки (алгоритм Томаса)} за $O(N)$.
  \item \textbf{Линейная элиминация на деревьях за $O(N)$:} пусть переходы происходят только по рёбрам дерева с корнем $r$ (между вершиной $u$, её родителем $p_u$ и сыновьями $v \in \operatorname{ch}(u)$). В обходе снизу вверх (DFS post-order) величина $E[u]$ для каждого поддерева параметризуется линейной функцией от значения в родителе:
  \[
    E[u] = k_u \cdot E[p_u] + b_u.
  \]
  Для терминальных вершин $u \in T$ полагаем $k_u = 0$, $b_u = 0$. Для $u \notin T$ подстановка $E[v] = k_v E[u] + b_v$ по всем $v \in \operatorname{ch}(u)$ в уравнение $E[u] = \bar{c}_u + p_{u, p_u} E[p_u] + \sum_{v \in \operatorname{ch}(u)} p_{u,v} E[v]$ даёт:
  \[
    \Delta_u = 1 - \sum_{v \in \operatorname{ch}(u)} p_{u,v} k_v, \qquad
    k_u = \frac{p_{u, p_u}}{\Delta_u}, \qquad
    b_u = \frac{\bar{c}_u + \sum_{v \in \operatorname{ch}(u)} p_{u,v} b_v}{\Delta_u}.
  \]
  В корне $p_{r, p_r} = 0$, поэтому $E[r] = b_r$; остальные значения при необходимости восстанавливаются проходом сверху вниз по формуле $E[u] = k_u E[p_u] + b_u$ за общее время $O(N)$.
\end{itemize}

% ------------------------------------------------------------------------------
\subsection{Лемма Бёрнсайда и теорема Пойа}

\textbf{Действие группы, орбиты и неподвижные точки.} Пусть конечная группа симметрий $(G, \cdot, e)$ действует на конечном множестве объектов $X$: каждому $g \in G$ и $x \in X$ сопоставлен элемент $g \cdot x \in X$, причём $e \cdot x = x$ и $(gh) \cdot x = g \cdot (h \cdot x)$.
\begin{itemize}
  \item \textbf{Орбита} элемента $x$: $G \cdot x = \{g \cdot x : g \in G\}$ (множество всех орбит обозначается $X / G$ — классы эквивалентности объектов с точностью до симметрий из $G$).
  \item \textbf{Стабилизатор} элемента $x$: $\Stab_G(x) = \{g \in G : g \cdot x = x\} \le G$; выполнено $|G \cdot x| \cdot |\Stab_G(x)| = |G|$.
  \item \textbf{Неподвижные точки} симметрии $g \in G$: $\Fix(g) = \{x \in X : g \cdot x = x\}$.
\end{itemize}

\textbf{Лемма Бёрнсайда.} Число различных орбит равно среднему числу неподвижных точек по всем элементам группы:
\[
  |X / G| = \frac{1}{|G|} \sum_{g \in G} |\Fix(g)|.
\]

\textbf{Теорема Редфилда—Пойа.} Пусть группа $G$ действует перестановками на множестве из $n = |D|$ позиций, и рассматривается множество $X = C^D$ всех раскрасок позиций в $m = |C|$ цветов. Пусть подстановка позиций, индуцированная элементом $g \in G$, разлагается на $c(g)$ независимых циклов, среди которых $c_j(g)$ циклов длины $j$ ($\sum_{j=1}^n j c_j(g) = n$, $\sum_{j=1}^n c_j(g) = c(g)$):
\begin{itemize}
  \item \textbf{Невесомая форма:} раскраска неподвижна относительно $g$ тогда и только тогда, когда все позиции внутри каждого цикла имеют один цвет, поэтому $|\Fix(g)| = m^{c(g)}$ и
  \[
    |C^D / G| = \frac{1}{|G|} \sum_{g \in G} m^{c(g)}.
  \]
  \item \textbf{Цикловой индекс и весовая форма (с ограничением на количества цветов):} цикловым индексом действия $G$ называется многочлен $Z_G(t_1, \dots, t_n) = \frac{1}{|G|} \sum_{g \in G} \prod_{j=1}^n t_j^{c_j(g)}$. Если цвету $r \in \{1, \dots, m\}$ сопоставлена формальная переменная $w_r$, то производящая функция неэквивалентных раскрасок равна:
  \[
    F(w_1, \dots, w_m) = Z_G\!\left(\sum_{r=1}^{m} w_r,\; \sum_{r=1}^{m} w_r^2,\; \dots,\; \sum_{r=1}^{m} w_r^n\right),
  \]
  где коэффициент при $w_1^{k_1} \dots w_m^{k_m}$ ($\sum_{r=1}^m k_r = n$) равен числу неэквивалентных раскрасок, содержащих ровно $k_r$ элементов цвета $r$.
\end{itemize}

\textbf{Канонические группы симметрий на окружности ($n$ вершин, $m$ цветов):}
\begin{itemize}
  \item \textbf{Ожерелья (циклическая группа поворотов $C_n$, $|C_n| = n$):} поворот на $k \in [0, n - 1]$ позиций имеет $c(\mathrm{rot}_k) = \gcd(n, k)$ циклов длины $n / \gcd(n, k)$. Группировка по делителям $d = \gcd(n, k) \mid n$:
  \[
    N_{\mathrm{rot}}(n, m) = \frac{1}{n} \sum_{k=0}^{n-1} m^{\gcd(n, k)} = \frac{1}{n} \sum_{d \mid n} \varphi\!\left(\frac{n}{d}\right) m^d.
  \]
  \item \textbf{Браслеты (диэдральная группа поворотов и отражений $D_n$, $|D_n| = 2n$):} включает $n$ поворотов и $n$ осевых отражений:
  \[
    N_{\mathrm{dih}}(n, m) =
    \begin{cases}
      \dfrac{1}{2n}\!\left(\displaystyle\sum_{d \mid n} \varphi\!\left(\frac{n}{d}\right) m^d + n \cdot m^{(n + 1)/2}\right), & n \text{ нечётно}, \\[10pt]
      \dfrac{1}{2n}\!\left(\displaystyle\sum_{d \mid n} \varphi\!\left(\frac{n}{d}\right) m^d + \frac{n}{2} m^{n/2 + 1} + \frac{n}{2} m^{n/2}\right), & n \text{ чётно.}
    \end{cases}
  \]
\end{itemize}

% ------------------------------------------------------------------------------
\subsection{Комбинаторика на графах}

\textbf{Код Прюфера.} Взаимно однозначное соответствие между множеством всех помеченных деревьев на вершинах $[n]$ ($n \ge 2$) и множеством кортежей $(p_1, \dots, p_{n - 2}) \in [n]^{n - 2}$:
\begin{itemize}
  \item \textbf{Кодирование $T \mapsto (p_1, \dots, p_{n - 2})$:} на шаге $i = 1, \dots, n - 2$ выбирается лист $v$ с минимальным номером, в $p_i$ записывается его единственный сосед, и лист $v$ удаляется из дерева.
  \item \textbf{Декодирование $(p_1, \dots, p_{n - 2}) \mapsto T$:} степень каждой вершины $v \in [n]$ в дереве равна $\deg(v) = 1 + |\{i : p_i = v\}|$. На шаге $i = 1, \dots, n - 2$ выбирается минимальная вершина $v$ с текущей степенью $1$, проводится ребро $(v, p_i)$ и обе степени уменьшаются на $1$; в конце ребром соединяются две оставшиеся вершины со степенью $1$. Оба преобразования реализуются за $O(n)$ с линейным указателем на минимальный лист.
\end{itemize}

\textbf{Формула Кэли и её обобщения:}
\begin{itemize}
  \item \textbf{Формула Кэли:} число помеченных деревьев на $n \ge 1$ вершинах равно $n^{n - 2}$ (при $n = 1$ равно $1$).
  \item \textbf{Деревья с заданными степенями:} при $n \ge 2$ число помеченных деревьев со степенями вершин $d_1, \dots, d_n \ge 1$ ($\sum_{i=1}^n d_i = 2n - 2$) равно $\frac{(n - 2)!}{(d_1 - 1)!\,(d_2 - 1)!\dots(d_n - 1)!}$.
  \item \textbf{Соединение компонент связности:} если граф на $n$ вершинах состоит из $k \ge 1$ связных компонент размеров $s_1, \dots, s_k$ ($\sum_{i=1}^k s_i = n$), то число способов добавить $k - 1$ рёбер до связного графа равно $n^{k - 2} \prod_{i=1}^k s_i$ (при $k = 1$ равно $1$). Число остовных лесов на $[n]$ из $m$ корневых деревьев с фиксированными корнями $\{1, \dots, m\}$ равно $m \cdot n^{n - m - 1}$.
\end{itemize}

\textbf{Матричная теорема Кирхгофа (Matrix-Tree Theorem).} Пусть $G = (V, E)$ — неориентированный мультиграф на вершинах $[n]$ без петель (петли не входят в остовные деревья и отбрасываются), где $a_{i,j} = a_{j,i}$ — кратность или сумма весов рёбер между $i$ и $j$ ($i \ne j$). \textbf{Матрица Лапласа (Кирхгофа)} $L = D - A \in \R^{n \times n}$ определяется как:
\[
  L_{i,i} = \sum_{k \ne i} a_{i,k}, \qquad L_{i,j} = -a_{i,j} \;\; (i \ne j).
\]
Для любого $r \in [n]$ главный минор порядка $(n - 1)$, полученный вычёркиванием $r$-й строки и $r$-го столбца матрицы $L$, равен числу остовных деревьев (а во взвешенном случае — сумме произведений весов рёбер по всем остовным деревьям $T$):
\[
  \tau(G) = \sum_{T \text{ — остов}} \;\prod_{e \in T} w(e) = \det L_{\bar{r}, \bar{r}} \qquad (\text{вычисляется методом Гаусса за } O(n^3)).
\]
\emph{Ориентированная версия (теорема Татта):} для орграфа с весами дуг $w_{u \to v}$ число корневых остовных деревьев с корнем $r$, ориентированных \textbf{от} корня $r$ ко всем вершинам (\emph{out-arborescence}), равно $\det L^{\mathrm{in}}_{\bar{r}, \bar{r}}$, где $L^{\mathrm{in}}_{v,v} = \sum_{u \ne v} w_{u \to v}$ и $L^{\mathrm{in}}_{u,v} = -w_{u \to v}$ ($u \ne v$). Для ориентации \textbf{к} корню $r$ (\emph{in-arborescence}) используется матрица исходящих степеней $L^{\mathrm{out}}_{u,u} = \sum_{v \ne u} w_{u \to v}$, $L^{\mathrm{out}}_{u,v} = -w_{u \to v}$.

% ==============================================================================
\section{Уровень 4: Продвинутый арсенал (CF 2300+)}
% ==============================================================================

\subsection{Инверсия Мёбиуса и биномиальная инверсия}

\textbf{Функция Мёбиуса $\mu : \N \to \{-1, 0, 1\}$.} Для канонического разложения $n = p_1^{\alpha_1} \dots p_r^{\alpha_r}$ ($p_i$ различны):
\[
  \mu(n) =
  \begin{cases}
    1, & n = 1, \\
    (-1)^r, & \alpha_1 = \dots = \alpha_r = 1 \text{ ($n$ свободно от квадратов)}, \\
    0, & \exists i : \alpha_i \ge 2.
  \end{cases}
\]
Функция $\mu$ мультипликативна ($\mu(ab) = \mu(a)\mu(b)$ при $\gcd(a, b) = 1$), вычисляется на отрезке $[1, N]$ линейным решетом за $O(N)$ и удовлетворяет фильтрующему тождеству:
\[
  \sum_{d \mid n} \mu(d) = \iverson{n = 1}, \qquad \sum_{d \mid n} \mu(d)\frac{n}{d} = \varphi(n).
\]

\textbf{Формулы инверсии Мёбиуса (по делителям и по кратным).} Для функций со значениями в абелевой группе:
\[
  g(n) = \sum_{d \mid n} f(d) \;\; (\forall n \ge 1) \quad \Longleftrightarrow \quad
  f(n) = \sum_{d \mid n} \mu(d)\, g\!\left(\frac{n}{d}\right) = \sum_{d \mid n} \mu\!\left(\frac{n}{d}\right) g(d);
\]
\[
  G(d) = \sum_{k=1}^{\lfloor M / d \rfloor} F(kd) \;\; (1 \le d \le M) \quad \Longleftrightarrow \quad
  F(d) = \sum_{k=1}^{\lfloor M / d \rfloor} \mu(k)\, G(kd).
\]

\textbf{Применение к задачам на массивы и графы с заданным $\gcd$:}
\begin{itemize}
  \item \textbf{Расщепление условия взаимной простоты:} подстановка $\iverson{\gcd(x, y) = 1} = \sum_{d \mid x,\, d \mid y} \mu(d)$ переводит условие на НОД в независимую делимость на $d$:
  \[
    \sum_{i, j} \iverson{\gcd(a_i, a_j) = g} w(i, j) = \sum_{k=1}^{\lfloor M / g \rfloor} \mu(k) \sum_{\substack{i : gk \mid a_i \\ j : gk \mid a_j}} w(i, j).
  \]
  \item \textbf{Пересчёт по кратным за $O(M \log M)$:} если $G(d)$ — ответ для объектов, у которых все элементы делятся на $d$ ($d \mid \gcd$), а $F(d)$ — ответ для объектов с $\gcd = d$, то все $F(1 \dots M)$ находятся либо через $F(d) = \sum_{k=1}^{\lfloor M/d \rfloor} \mu(k) G(kd)$, либо обратным вычитанием от $d = M$ вниз до $1$: $F(d) = G(d) - \sum_{k=2}^{\lfloor M/d \rfloor} F(kd)$ в силу гармонической оценки $\sum_{d=1}^M \lfloor M / d \rfloor = O(M \log M)$.
  \item \textbf{Подграфы кратных весов:} для подсчёта путей, компонент или остовных деревьев с заданным $\gcd$ весов рёбер для каждого $d$ строится подграф $H_d$ только из рёбер с $d \mid w_e$. Суммарное число рёбер во всех непустых $H_d$ равно $\sum_{e \in E} \tau(w_e)$, где $\tau(x)$ — число делителей ($\max_{x \le 10^6} \tau(x) = 240$).
\end{itemize}

\textbf{Биномиальная инверсия (переход между «выбрано $k$ условий», «ровно $k$» и «хотя бы $k$»).}
\begin{itemize}
  \item \textbf{Нижнетреугольная (классическая) форма:}
  \[
    b_n = \sum_{k=0}^{n} \binom{n}{k} a_k \quad \Longleftrightarrow \quad a_n = \sum_{k=0}^{n} (-1)^{n - k} \binom{n}{k} b_k.
  \]
  \item \textbf{Верхнетреугольная форма (обобщённая ФВИ для ДП):} пусть заданы $n$ свойств $A_1, \dots, A_n \subseteq U$; обозначим через $E_k$ число объектов, обладающих \textbf{ровно} $k$ свойствами, а через $G_k = \sum_{S \subseteq [n],\, |S| = k} \left|\bigcap_{i \in S} A_i\right|$ — сумму размеров пересечений по всем наборам из $k$ фиксированных свойств (каждый объект с ровно $j \ge k$ свойствами учитывается в $G_k$ с кратностью $\binom{j}{k}$). Тогда для всех $0 \le k \le n$:
  \[
    G_k = \sum_{j=k}^{n} \binom{j}{k} E_j \quad \Longleftrightarrow \quad
    E_k = \sum_{j=k}^{n} (-1)^{j - k} \binom{j}{k} G_j.
  \]
  Число объектов, обладающих \textbf{хотя бы} $k$ свойствами без кратного пересчёта ($L_k = \sum_{j=k}^n E_j$, $k \ge 1$), выражается напрямую через $G_j$:
  \[
    L_k = \sum_{j=k}^{n} (-1)^{j - k} \binom{j - 1}{k - 1} G_j.
  \]
  Весь вектор $(E_0, \dots, E_n)$ вычисляется по $(G_0, \dots, G_n)$ за $O(n^2)$ по формуле или за $O(n \log n)$ одной свёрткой полиномов, так как $k! E_k = \sum_{u=0}^{n - k} \frac{(-1)^u}{u!} \cdot (k + u)! G_{k + u}$.
\end{itemize}

% ------------------------------------------------------------------------------
\subsection{Производящие функции (OGF и EGF) + NTT}

\textbf{Обыкновенные производящие функции (OGF).} Последовательности $(a_n)_{n \ge 0}$ над полем $\K$ (обычно число \emph{непомеченных} объектов веса $n$) сопоставляется формальный степенной ряд $A(x) = \sum_{n=0}^\infty a_n x^n \in \K[[x]]$:
\begin{itemize}
  \item \textbf{Операции:} $[x^n](A(x) + B(x)) = a_n + b_n$; $[x^n](A(x)B(x)) = \sum_{k=0}^n a_k b_{n - k}$ (свёртка Коши — упорядоченные пары объектов суммарного веса $n$); $[x^n]\frac{A(x)}{1 - x} = \sum_{k=0}^n a_k$; при $a_0 = 0$ ряд $\operatorname{SEQ}(A) = \frac{1}{1 - A(x)} = \sum_{k=0}^\infty A(x)^k$ перечисляет конечные последовательности непустых объектов из $A$.
  \item \textbf{Базовые табличные OGF} (обобщённый биномиальный коэффициент $\binom{\alpha}{n} = \frac{\fall{\alpha}{n}}{n!}$ для любого $\alpha \in \K$):
  \[
    \frac{1}{1 - x} = \sum_{n=0}^{\infty} x^n, \quad
    (1 + x)^\alpha = \sum_{n=0}^{\infty} \binom{\alpha}{n} x^n, \quad
    \frac{1}{(1 - x)^k} = \sum_{n=0}^{\infty} \binom{n + k - 1}{k - 1} x^n, \quad
    \sum_{n=0}^{\infty} C_n x^n = \frac{1 - \sqrt{1 - 4x}}{2x}.
  \]
\end{itemize}

\textbf{Экспоненциальные производящие функции (EGF).} Последовательности $(a_n)_{n \ge 0}$ (число \emph{помеченных} структур на множестве из $n$ меток) сопоставляется ряд $\widehat{A}(x) = \sum_{n=0}^\infty a_n \frac{x^n}{n!}$, так что $a_n = n! [x^n]\widehat{A}(x)$:
\begin{itemize}
  \item \textbf{Биномиальная свёртка:} $\widehat{C}(x) = \widehat{A}(x)\widehat{B}(x) \iff c_n = \sum_{k=0}^n \binom{n}{k} a_k b_{n - k}$ (упорядоченные разбиения множества меток $[n] = S_1 \sqcup S_2$ со структурами из $\mathcal{A}$ на $S_1$ и из $\mathcal{B}$ на $S_2$). Производная сдвигает индекс: $\widehat{A}'(x) = \sum_{n=0}^\infty a_{n+1}\frac{x^n}{n!}$.
  \item \textbf{Базовые EGF:} $e^x \leftrightarrow a_n = 1$ (множества); $\frac{1}{1 - x} \leftrightarrow n!$ (перестановки); $\ln\frac{1}{1 - x} \leftrightarrow (n - 1)!$ ($n \ge 1$, ориентированные циклы); $\frac{e^{-x}}{1 - x} \leftrightarrow !n$ (беспорядки); $\frac{(e^x - 1)^k}{k!} \leftrightarrow \stirlingsecond{n}{k}$ (разбиения на $k$ блоков).
  \item \textbf{Экспоненциальная формула (связь связных и всех помеченных объектов):} пусть $\widehat{C}(x) = \sum_{n=1}^\infty c_n \frac{x^n}{n!}$ ($c_0 = 0$) — EGF непустых \emph{связных} помеченных структур, а $\widehat{G}(x) = \sum_{n=0}^\infty g_n \frac{x^n}{n!}$ ($g_0 = 1$) — EGF объектов, представляющих собой произвольное неупорядоченное множество связных компонент из $\mathcal{C}$. Тогда:
  \[
    \widehat{G}(x) = \exp\!\bigl(\widehat{C}(x)\bigr) \quad \Longleftrightarrow \quad \widehat{C}(x) = \ln \widehat{G}(x), \qquad
    \text{EGF для ровно } k \text{ компонент} = \frac{\widehat{C}(x)^k}{k!}.
  \]
  \emph{Пример (связные простые помеченные графы):} $g_n = 2^{\binom{n}{2}}$, поэтому числа связных графов $c_1, \dots, c_N$ находятся взятием полиномиального логарифма $\widehat{C}(x) = \ln\!\bigl(\sum_{n=0}^N 2^{\binom{n}{2}} \frac{x^n}{n!}\bigr) \bmod x^{N+1}$ за $O(N \log N)$.
\end{itemize}

\textbf{NTT (Number-Theoretic Transform) по модулю $p = 998244353$.} Для простого $p = c \cdot 2^K + 1$ с первообразным корнем $g$ мультипликативной группы $\F_p^\times$ ($p = 998244353 = 119 \cdot 2^{23} + 1$, $g = 3$) для любой степени двойки $N = 2^m \le 2^K$ существует примитивный корень степени $N$: $\omega_N = g^{(p - 1)/N} \bmod p$ ($\omega_N^N \equiv 1$, $\omega_N^{N/2} \equiv p - 1 \pmod p$). Прямое и обратное преобразования вектора длины $N$:
\[
  \hat{a}_j = \sum_{k=0}^{N - 1} a_k \omega_N^{jk} \bmod p, \qquad
  a_k = N^{-1} \sum_{j=0}^{N - 1} \hat{a}_j \omega_N^{-jk} \bmod p \qquad (0 \le j, k < N).
\]
Поразрядный алгоритм Кули—Тьюки вычисляет оба преобразования за $O(N \log N)$, что даёт перемножение многочленов степеней $< n$ и $< m$ за $O(N \log N)$ при $N = 2^{\lceil \log_2(n + m - 1) \rceil}$.

\textbf{Операции $\inv, \ln, \exp$ над формальными рядами по модулю $x^N$ за $O(N \log N)$} (удвоение точности $m = 1, 2, 4, \dots$ методом Ньютона над $\F_p[[x]]$, при $N \le p$):
\begin{enumerate}
  \item \textbf{Обратный ряд $B(x) = A(x)^{-1} \bmod x^N$} (определён при $a_0 = [x^0]A(x) \not\equiv 0 \pmod p$):
  \[
    B_1(x) = a_0^{-1} \bmod p, \qquad
    B_{2m}(x) \equiv B_m(x)\bigl(2 - A(x)B_m(x)\bigr) \pmod{x^{2m}}.
  \]
  \item \textbf{Логарифм $\ln A(x) \bmod x^N$} (определён при $[x^0]A(x) = 1$, причём $[x^0]\ln A(x) = 0$; формальный интеграл задаётся линейно по правилу $\int x^k\,dx = \frac{x^{k + 1}}{k + 1}$):
  \[
    \ln A(x) \equiv \int \Bigl(A'(x) \cdot A(x)^{-1} \bmod x^{N - 1}\Bigr) dx \pmod{x^N}.
  \]
  \item \textbf{Экспонента $\exp A(x) \bmod x^N$} (определена при $[x^0]A(x) = 0$, с $[x^0]\exp A(x) = 1$):
  \[
    E_1(x) = 1, \qquad
    E_{2m}(x) \equiv E_m(x)\bigl(1 - \ln E_m(x) + A(x)\bigr) \pmod{x^{2m}}.
  \]
  \item \textbf{Степень ряда:} если $A(x) = c x^d \tilde{A}(x)$ с $[x^0]\tilde{A}(x) = 1$ и $c \not\equiv 0$, то $A(x)^k \equiv c^k x^{kd} \exp(k \ln \tilde{A}(x)) \pmod{x^N}$.
\end{enumerate}

% ------------------------------------------------------------------------------
\subsection{Специальные приёмы}

\textbf{Max-Min Inclusion-Exclusion (тождество минимума и максимума).} Для любого непустого конечного множества вещественных чисел $S$ и для любого набора случайных величин $\{X_i\}_{i \in S}$ с конечными матожиданиями:
\[
  \max(S) = \sum_{\emptyset \ne T \subseteq S} (-1)^{|T| - 1} \min(T), \qquad
  \E\!\left[\max_{i \in S} X_i\right] = \sum_{\emptyset \ne T \subseteq S} (-1)^{|T| - 1} \E\!\left[\min_{i \in T} X_i\right].
\]
Если $X_i$ — момент первого появления элемента $i$ в серии независимых испытаний, где элемент $i$ выпадает с вероятностью $p_i$, то $\max_{i \in S} X_i$ — время сбора всего множества $S$, а $\min_{i \in T} X_i \sim \Geom\bigl(\sum_{i \in T} p_i\bigr)$, откуда $\E[\min_{i \in T} X_i] = \bigl(\sum_{i \in T} p_i\bigr)^{-1}$.
Обобщение для $k$-го по величине максимума $\max^{(k)}(S)$ ($1 \le k \le |S|$, где $\max^{(1)} = \max$):
\[
  \E\!\left[\max\nolimits^{(k)}_{i \in S} X_i\right] = \sum_{\substack{T \subseteq S \\ |T| \ge k}} (-1)^{|T| - k} \binom{|T| - 1}{k - 1} \E\!\left[\min_{i \in T} X_i\right].
\]

\textbf{Лемма Линдстрёма—Гесселя—Вьенно (LGV — непересекающиеся пути в DAG).} Пусть $G = (V, E)$ — конечный ориентированный ациклический граф с весами дуг $w : E \to R$ ($R$ — коммутативное кольцо; при простом подсчёте путей $w(e) = 1$). Вес пути $P$ равен $w(P) = \prod_{e \in P} w(e)$, а суммарный вес всех путей из $u$ в $v$ равен $e(u, v) = \sum_{P : u \leadsto v} w(P)$. Для фиксированных наборов из $k$ стартовых вершин $A = (a_1, \dots, a_k)$ и $k$ финишных вершин $B = (b_1, \dots, b_k)$ рассмотрим матрицу $M = \bigl(e(a_i, b_j)\bigr)_{i,j=1}^k$. Тогда:
\[
  \det M = \sum_{\sigma \in S_k} \sgn(\sigma) \sum_{\substack{\mathcal{P} = (P_1, \dots, P_k) \\ P_i : a_i \leadsto b_{\sigma(i)},\; V(P_i) \cap V(P_j) = \emptyset}} \prod_{i=1}^{k} w(P_i),
\]
где $\sgn(\sigma) = (-1)^{\inv(\sigma)}$, а внутренняя сумма берётся только по наборам попарно \textbf{вершинно-непересекающихся} путей.

\emph{Следствие для решёток и планарных DAG:} если по геометрии графа любой вершинно-непересекающийся набор путей $a_i \leadsto b_{\sigma(i)}$ возможен \textbf{только при тождественной перестановке} $\sigma = \mathrm{id}$ (например, монотонные пути на $\Z^2$ с упорядоченными по границе источниками и стоками), то $\det M$ в точности равен числу (суммарному весу) систем из $k$ вершинно-непересекающихся путей $P_i : a_i \leadsto b_i$. Сложность: $O(k(|V| + |E|) + k^3)$.

\textbf{Вероятностный метод и рандомизация:}
\begin{itemize}
  \item \textbf{Базовые принципы вероятностного метода:}
  \begin{enumerate}
    \item \emph{Аддитивная оценка (Union Bound):} $P\bigl(\bigcup_{i=1}^m B_i\bigr) \le \sum_{i=1}^m P(B_i)$; если $\sum_{i=1}^m P(B_i) < 1$, то существует исход $\omega \in \Omega$, свободный от всех «плохих» свойств $B_i$.
    \item \emph{Метод первого момента:} для любой СВ $X$ существуют $\omega_1, \omega_2 \in \Omega$, на которых $X(\omega_1) \ge \E[X]$ и $X(\omega_2) \le \E[X]$.
    \item \emph{Лемма Шварца—Зиппеля:} если $P(x_1, \dots, x_n) \not\equiv 0$ — ненулевой многочлен суммарной степени $\deg P \le d$ над полем $\F$ и точки $r_1, \dots, r_n$ выбираются независимо и равновероятно из конечного множества $S \subseteq \F$, то $P\bigl(P(r_1, \dots, r_n) = 0\bigr) \le \frac{d}{|S|}$.
  \end{enumerate}

  \item \textbf{Парадокс дней рождения:} при независимом равновероятном выборе $k$ элементов из множества размера $M$ вероятность попарной различности оценивается сверху как:
  \[
    P(\text{нет коллизий}) = \prod_{i=0}^{k - 1}\left(1 - \frac{i}{M}\right) \le \exp\!\left(-\frac{k(k - 1)}{2M}\right), \qquad
    P(\text{есть коллизия}) \le \frac{k(k - 1)}{2M}.
  \]
  Коллизия возникает с вероятностью $\ge 1/2$ при $k \approx 1{,}177\sqrt{M}$ ($\Theta(\sqrt{M})$). При попарном сравнении или хранении в таблице $k$ случайных хешей для малой вероятности ошибки необходимо выбирать размер пространства хешей $M \gg k^2$ (при $k = 10^6$ требуется $M \sim 2^{64}$ или пара независимых $32$-битных модулей).

  \item \textbf{Zobrist-хеширование:} каждому базовому элементу (или паре «позиция—состояние») $x \in \mathcal{U}$ сопоставляется независимое случайное число $h(x) \sim \Unif(\{0, \dots, 2^{64} - 1\})$:
  \begin{itemize}
    \item \emph{Множества и чётность вхождений:} $H_\oplus(S) = \bigoplus_{x \in S} h(x)$. Операция не зависит от порядка, обновляется при вставке/удалении за $O(1)$ ($H_\oplus(A \triangle B) = H_\oplus(A) \oplus H_\oplus(B)$); для любого непустого множества с нечётными кратностями $P(H_\oplus = 0) = 2^{-64}$.
    \item \emph{Мультимножества:} $H_+(M) = \sum_{x \in M} h(x) \bmod 2^{64}$ (либо по большому простому модулю $P$); для проверки кратности вхождений делителю $K$ каждая кратность поддерживается по модулю $K$ с весом $h(x)$ (или суммой случайных векторов над $\Z_K$).
  \end{itemize}

  \item \textbf{Алгоритм Фрейвалдса (проверка матричного произведения $A \cdot B = C$ за $O(k n^2)$):} для квадратных матриц $A, B, C$ размера $n \times n$ на каждой из $k$ независимых итераций генерируется случайный вектор $r \in \{0, 1\}^n$ ($r_i \sim \Unif(\{0, 1\})$) и проверяется равенство векторов длины $n$:
  \[
    A(Br) = Cr \qquad \bigl(\text{три умножения матрицы на вектор за } O(n^2)\bigr).
  \]
  Если $AB = C$, равенство выполняется всегда; если $AB \ne C$, вероятность ложного совпадения на одной итерации не превосходит $1/2$, а за $k$ независимых итераций — не превосходит $2^{-k}$ (при $k = 40$ ошибка $\le 2^{-40} < 10^{-12}$).
\end{itemize}

\end{document}
